\documentclass[11pt,a4paper]{ctexart}
\usepackage{amsmath,amssymb,amsthm}
\usepackage{enumitem}
\usepackage[margin=2.4cm]{geometry}
\usepackage[colorlinks=true,linkcolor=blue,citecolor=blue,urlcolor=blue]{hyperref}

\theoremstyle{definition}
\newtheorem{definition}{定义}[section]

\theoremstyle{plain}
\newtheorem{proposition}[definition]{命题}
\newtheorem{lemma}[definition]{引理}
\newtheorem{theorem}[definition]{定理}
\newtheorem{corollary}[definition]{推论}

\renewcommand{\proofname}{证明}

\newcommand{\Stab}{\operatorname{Stab}}
\newcommand{\im}{\operatorname{im}}
\newcommand{\cof}{\operatorname{cof}}
\newcommand{\Reg}{\operatorname{Reg}}
\newcommand{\Ord}{\operatorname{Ord}}
\newcommand{\id}{\mathrm{id}}
\newcommand{\Set}{\mathbf{Set}}

\title{集合函子 $G_{\kappa,\theta}$ 精确稳定阶 $\Stab(G_{\kappa,\theta}) = \kappa + \theta$ 的证明}
\date{}

\begin{document}
\maketitle

\section{函子定义、终端序列与基本性质}

设 $\kappa, \theta$ 为无穷正则基数，且满足 $\omega \le \theta \le \kappa$。记 $\Reg_{<\theta}$ 为所有严格小于 $\theta$ 的无穷正则基数 $\mu$ 构成的集合。对每个 $\mu \in \Reg_{<\theta}$，在函数空间 $X^\mu$ 上定义等价关系：
\[
  a \equiv_{<\mu} b \iff \bigl|\{ i < \mu : a(i) \neq b(i) \}\bigr| < \mu.
\]
由于 $\mu$ 为正则基数（从而两个小差异集之并的基数仍严格小于 $\mu$），$\equiv_{<\mu}$ 是等价关系。据此定义集合端函子 $G_{\kappa,\theta} : \Set \to \Set$：
\[
  G_{\kappa,\theta}(X) = \mathcal{P}_{<\kappa}(X) \;\sqcup\; \bigsqcup_{\mu \in \Reg_{<\theta}} \bigl(X^\mu / {\equiv_{<\mu}}\bigr),
\]
其中 $\mathcal{P}_{<\kappa}(X) = \{ S \subseteq X : |S| < \kappa \}$。给定映射 $h : X \to Y$，$G_{\kappa,\theta}(h)$ 在左分支取直接像 $S \mapsto h[S]$（由 $|h[S]| \le |S| < \kappa$ 良定义），在右分支逐分量复合 $[a] \mapsto [h \circ a]$（由 $\{ i : h(a(i)) \neq h(b(i)) \} \subseteq \{ i : a(i) \neq b(i) \}$ 良定义）。为记号简洁，以下亦将左分支嵌入记作 $\mathsf{L}(S)$、右分支嵌入记作 $\mathsf{R}(\mu, [a])$。

\subsection*{终端序列 $(A_\alpha, f_\alpha^\beta)$}

在序数类 $\Ord$ 上通过超限递归定义 $G_{\kappa,\theta}$ 的终端序列（terminal sequence）：
\begin{itemize}[leftmargin=2em]
  \item \textbf{初项}：$A_0 = 1 = \{*\}$（终对象）；
  \item \textbf{后继步}：$A_{\alpha+1} = G_{\kappa,\theta}(A_\alpha)$，衔接态射满足 $f_{\alpha+1}^{\beta+1} = G_{\kappa,\theta}(f_\alpha^\beta)$；
  \item \textbf{极限步}：对极限序数 $\lambda$，
    \[
      A_\lambda = \varprojlim_{\alpha < \lambda} A_\alpha = \Bigl\{ (x_\alpha)_{\alpha < \lambda} \in \prod_{\alpha<\lambda} A_\alpha \;\Big|\; \forall \alpha \le \beta < \lambda,\; f_\alpha^\beta(x_\beta) = x_\alpha \Bigr\},
    \]
    且 $f_\alpha^\lambda$ 为第 $\alpha$ 个投影映射。
\end{itemize}
全体过渡映射满足 $f_\alpha^\alpha = \id$ 与 $f_\alpha^\gamma = f_\alpha^\beta \circ f_\beta^\gamma$（$\alpha \le \beta \le \gamma$）。此外，因 $\emptyset \in \mathcal{P}_{<\kappa}(1)$，恒有 $A_1 \neq \emptyset$，从而对一切序数 $\alpha$ 均有 $A_\alpha \neq \emptyset$。若 $f_\gamma^{\gamma+1}$ 为双射，则称 $G_{\kappa,\theta}$ 在 $\gamma$ 处\textbf{稳定}（此时对一切 $\delta \ge \gamma$，$f_\gamma^\delta$ 与 $f_\delta^{\delta+1}$ 均为双射）；使之稳定的最小序数称为精确稳定阶 $\Stab(G_{\kappa,\theta})$。

\begin{proposition}[$\kappa$-可及性与保持单射]
当 $\kappa$ 为无穷正则基数且 $\theta \le \kappa$ 时：
\begin{enumerate}[label=(\arabic*)]
  \item $G_{\kappa,\theta}$ \textbf{保持单射}：若 $h : X \hookrightarrow Y$ 为单射，则 $G_{\kappa,\theta}(h)$ 亦为单射；
  \item $G_{\kappa,\theta}$ \textbf{是 $\kappa$-可及的}：对任意集合 $X$ 与 $u \in G_{\kappa,\theta}(X)$，存在基数 $|Y| < \kappa$ 的子集 $Y \subseteq X$ 使 $u \in \im\bigl(G_{\kappa,\theta}(\iota_{Y \hookrightarrow X})\bigr)$。
\end{enumerate}
\end{proposition}

\begin{proof}
(1) 左右不交并分支互不相混。在左分支，若 $h[S_1] = h[S_2]$ 则由 $h$ 单射得 $S_1 = S_2$；在右分支，若 $[h \circ a] = [h \circ b]$，则 $\{ i : a(i) \neq b(i) \} = \{ i : h(a(i)) \neq h(b(i)) \}$ 的基数严格小于 $\mu$，故 $[a] = [b]$。

(2) 若 $u = \mathsf{L}(S)$，直接取 $Y = S$（$|S| < \kappa$）；若 $u = \mathsf{R}(\mu, [a])$（$\mu < \theta \le \kappa$），取 $Y = a[\mu] \subseteq X$，则 $|Y| \le \mu < \kappa$ 且 $a$ 通过 $Y$ 分解。
\end{proof}

\begin{proposition}[尾部单态性]
若 $\kappa$ 为无穷正则基数且 $\theta \le \kappa$，则对一切 $\kappa \le \alpha \le \beta$，转移映射 $f_\alpha^\beta : A_\beta \to A_\alpha$ 均为单射；特别地，$f_\kappa^{\kappa+1}$ 是单射。
\end{proposition}

\begin{proof}
先证 $f_\kappa^{\kappa+1}$ 单射。若 $Y \subseteq A_\kappa$ 满足 $|Y| < \kappa$，则对任意 $x \neq y \in Y$ 存在 $\alpha(x,y) < \kappa$ 使 $f_{\alpha(x,y)}^\kappa(x) \neq f_{\alpha(x,y)}^\kappa(y)$；因 $|Y \times Y| < \kappa$ 且 $\kappa$ 正则，$\alpha_0 := \sup_{x \neq y \in Y} \alpha(x,y) < \kappa$，于是 $f_{\alpha_0}^\kappa|_Y$ 为单射。

现设 $x, y \in A_{\kappa+1} = G_{\kappa,\theta}(A_\kappa)$ 且 $f_\kappa^{\kappa+1}(x) = f_\kappa^{\kappa+1}(y)$。由 $\kappa$-可及性取基数 $< \kappa$ 的支集 $Y_x, Y_y \subseteq A_\kappa$，并取 $Y = Y_x \cup Y_y \cup \{a_0\}$（$a_0 \in A_\kappa$ 任取），仍有 $0 < |Y| < \kappa$。取如上 $\alpha < \kappa$ 使 $f_\alpha^\kappa|_Y$ 单射，则 $f_\alpha^\kappa \circ \iota_Y : Y \hookrightarrow A_\alpha$ 为非空集合间的单射，必存在左逆 $r : A_\alpha \to Y$。由于
\[
  G_{\kappa,\theta}(f_\alpha^\kappa)(x) = f_\alpha^{\kappa+1}(x) = f_\alpha^\kappa\bigl(f_\kappa^{\kappa+1}(x)\bigr) = f_\alpha^\kappa\bigl(f_\kappa^{\kappa+1}(y)\bigr) = G_{\kappa,\theta}(f_\alpha^\kappa)(y),
\]
而 $x, y$ 分别由 $Y$ 上元素经 $\iota_Y$ 映入，在两端作用 $G_{\kappa,\theta}(r)$ 即得 $x = y$。

再对 $\beta \ge \kappa$ 作超限归纳：后继步由 $G_{\kappa,\theta}$ 保持单射传递，极限步由 $\varprojlim$ 投影唯一性传递，故全部 $f_\alpha^\beta$（$\kappa \le \alpha \le \beta$）皆为单射。
\end{proof}

\section{第一阶段：$f_\kappa^{\kappa+1}$ 不满射，从而 $\Stab(G_{\kappa,\theta}) \ge \kappa$}

本阶段仅用到左分支 $\mathcal{P}_{<\kappa}(X)$ 与 $\kappa$ 的正则性（对 $\theta$ 无任何额外要求）。

\subsection*{2.1 对角家族 $q_\alpha(\beta)$ 的构造}

对 $\alpha \le \kappa$ 与全部序数 $\beta \in \Ord$，按关于 $\alpha$ 的超限递归定义元素 $q_\alpha(\beta) \in A_\alpha$：
\begin{enumerate}[label=(\arabic*)]
  \item $\alpha = 0$：$q_0(\beta) := * \in A_0$；
  \item 后继步（$\alpha = \delta + 1 \le \kappa$）：
    \[
      q_{\delta+1}(\beta) := \mathsf{L}\Bigl(\bigl\{ q_\delta(\gamma) \mid \gamma < \min(\beta,\; \delta+1) \bigr\}\Bigr) \in G_{\kappa,\theta}(A_\delta) = A_{\delta+1}
    \]
    （因 $\delta+1 < \kappa$，该集合基数 $\le |\delta+1| < \kappa$，确实属于 $\mathcal{P}_{<\kappa}(A_\delta)$）；
  \item 极限步（$l \le \kappa$ 为极限序数）：$q_l(\beta) := \bigl(q_\alpha(\beta)\bigr)_{\alpha < l} \in A_l$。
\end{enumerate}
在严格构造中，极限元在 $A_l$ 中的相容性可与下面两条核心引理同步归纳证明。

\begin{lemma}[截断等同判据]
对一切 $\alpha \le \kappa$ 与一切序数 $\beta, \gamma \in \Ord$，
\[
  q_\alpha(\beta) = q_\alpha(\gamma) \iff \min(\alpha, \beta) = \min(\alpha, \gamma).
\]
\end{lemma}

\begin{proof}
对 $\alpha$ 作超限归纳：
\begin{itemize}[leftmargin=2em]
  \item \textbf{基始步}（$\alpha = 0$）：$q_0(\beta) = * = q_0(\gamma)$，且 $\min(0, \beta) = 0 = \min(0, \gamma)$，成立。
  \item \textbf{后继步}（$\alpha = \delta + 1$）：
    \[
      q_{\delta+1}(\beta) = q_{\delta+1}(\gamma) \iff \bigl\{q_\delta(i) \mid i < \min(\beta, \delta+1)\bigr\} = \bigl\{q_\delta(i) \mid i < \min(\gamma, \delta+1)\bigr\}.
    \]
    对一切 $i, j < \delta+1$（即 $i, j \le \delta$），由归纳假设有
    \[
      q_\delta(i) = q_\delta(j) \iff \min(\delta, i) = \min(\delta, j) \iff i = j.
    \]
    故映射 $i \mapsto q_\delta(i)$ 在 $\{i : i < \delta+1\}$ 上是单射。从而两个集合相等当且仅当指标上界相等，即 $\min(\beta, \delta+1) = \min(\gamma, \delta+1)$。
  \item \textbf{极限步}（$l$ 为极限序数）：
    \[
      q_l(\beta) = q_l(\gamma) \iff \forall \alpha < l,\; q_\alpha(\beta) = q_\alpha(\gamma) \iff \forall \alpha < l,\; \min(\alpha, \beta) = \min(\alpha, \gamma).
    \]
    \begin{itemize}
      \item[($\Leftarrow$)] 若 $\min(l, \beta) = \min(l, \gamma)$，则对一切 $\alpha < l$ 有 $\min(\alpha, \beta) = \min(\alpha, \min(l, \beta)) = \min(\alpha, \min(l, \gamma)) = \min(\alpha, \gamma)$。
      \item[($\Rightarrow$)] 设 $\forall \alpha < l,\; \min(\alpha, \beta) = \min(\alpha, \gamma)$，反设 $\min(l, \beta) \neq \min(l, \gamma)$。不妨设 $\mu := \min(l, \beta) < \min(l, \gamma) \le l$。因 $\mu < l$ 且 $l$ 为极限序数，有 $\mu + 1 < l$。取 $\alpha = \mu + 1$：因 $\mu < \min(l, \gamma) \le \gamma$ 故 $\min(\mu+1, \gamma) = \mu + 1$，而 $\beta = \mu < \mu + 1$ 故 $\min(\mu+1, \beta) = \mu$，与假设矛盾。 \qedhere
    \end{itemize}
\end{itemize}
\end{proof}

\begin{lemma}[投影相容性]
对一切 $\alpha \le \eta \le \kappa$ 与一切 $\beta \in \Ord$，均有 $f_\alpha^\eta\bigl(q_\eta(\beta)\bigr) = q_\alpha(\beta)$。
\end{lemma}

\begin{proof}
对 $\eta$ 作超限归纳。当 $\alpha = \eta$ 时由 $f_\eta^\eta = \id$ 平凡成立；下设 $\alpha < \eta$。
\begin{itemize}[leftmargin=2em]
  \item \textbf{后继步}（$\eta = \xi + 1$）：只需证 $f_\xi^{\xi+1}\bigl(q_{\xi+1}(\beta)\bigr) = q_\xi(\beta)$（再复合 $f_\alpha^\xi$ 与归纳假设即得对所有 $\alpha \le \xi$ 成立）。对 $\xi$ 分类讨论：
    \begin{enumerate}[label=(\roman*)]
      \item $\xi = 0$：$A_0 = \{*\}$，故 $f_0^1(q_1(\beta)) = * = q_0(\beta)$。
      \item $\xi = \delta + 1$（后继序数）：此时 $f_{\delta+1}^{\delta+2} = G_{\kappa,\theta}(f_\delta^{\delta+1})$，在左分支上即逐元取像：
        \[
          f_{\delta+1}^{\delta+2}\bigl(q_{\delta+2}(\beta)\bigr) = \mathsf{L}\Bigl(\bigl\{ f_\delta^{\delta+1}(q_{\delta+1}(i)) \mid i < \min(\beta, \delta+2) \bigr\}\Bigr) = \mathsf{L}\Bigl(\bigl\{ q_\delta(i) \mid i < \min(\beta, \delta+2) \bigr\}\Bigr).
        \]
        因由截断等同判据有 $q_\delta(\delta+1) = q_\delta(\delta)$，上式大括号内的集合正等于 $\{ q_\delta(i) \mid i < \min(\beta, \delta+1) \}$，故等于 $q_{\delta+1}(\beta)$。
      \item $\xi$ 为极限序数：只需比较二者到每个 $\alpha < \xi$ 的投影。因 $\xi$ 极限，对任意 $\alpha < \xi$ 有 $\alpha + 1 < \xi$；利用 $f_\alpha^\xi \circ f_\xi^{\xi+1} = f_\alpha^{\alpha+1} \circ G_{\kappa,\theta}(f_\alpha^\xi)$ 及归纳假设，得
        \[
          f_\alpha^\xi\bigl(f_\xi^{\xi+1}(q_{\xi+1}(\beta))\bigr) = f_\alpha^{\alpha+1}\Bigl(\mathsf{L}\bigl\{ f_\alpha^\xi(q_\xi(i)) \mid i < \min(\beta, \xi+1) \bigr\}\Bigr) = f_\alpha^{\alpha+1}\Bigl(\mathsf{L}\bigl\{ q_\alpha(i) \mid i < \min(\beta, \xi+1) \bigr\}\Bigr).
        \]
        由于 $\forall i \ge \alpha+1$ 均有 $q_\alpha(i) = q_\alpha(\alpha)$，集合 $\{ q_\alpha(i) \mid i < \min(\beta, \xi+1) \}$ 恰好等于 $\{ q_\alpha(i) \mid i < \min(\beta, \alpha+1) \}$，即其内层就是 $q_{\alpha+1}(\beta)$，故上式等于 $f_\alpha^{\alpha+1}(q_{\alpha+1}(\beta)) = q_\alpha(\beta)$。
    \end{enumerate}
  \item \textbf{极限步}（$\eta$ 为极限序数）：对一切 $\alpha < \eta$，$f_\alpha^\eta(q_\eta(\beta))$ 按极限投影定义即为第 $\alpha$ 分量 $q_\alpha(\beta)$。 \qedhere
\end{itemize}
\end{proof}

\begin{theorem}[第 $\kappa$ 级非满射性]
令 $b_\kappa := q_\kappa(\kappa) \in A_\kappa$。则 $b_\kappa \notin \im(f_\kappa^{\kappa+1})$；从而 $f_\kappa^{\kappa+1}$ 不满射，且对一切 $\delta < \kappa$，$G_{\kappa,\theta}$ 在 $\delta$ 处均不稳定。
\end{theorem}

\begin{proof}
反设存在 $x \in A_{\kappa+1} = G_{\kappa,\theta}(A_\kappa)$ 使 $f_\kappa^{\kappa+1}(x) = b_\kappa$。分两分支导出矛盾：
\begin{itemize}[leftmargin=2em]
  \item \textbf{右分支 $x = \mathsf{R}(\mu, [a])$}：将 $f_\kappa^{\kappa+1}(x) = q_\kappa(\kappa)$ 用 $f_1^\kappa$ 投影到第 $1$ 级。左端为 $f_1^{\kappa+1}(\mathsf{R}(\mu, [a])) = \mathsf{R}(\mu, [f_0^\kappa \circ a])$ 仍落在右分支；右端按投影相容性引理等于 $q_1(\kappa) = \mathsf{L}(\{*\})$ 落在左分支，矛盾。
  \item \textbf{左分支 $x = \mathsf{L}(S)$}（其中 $S \subseteq A_\kappa$ 且 $|S| < \kappa$）：
    \begin{enumerate}[label=(\arabic*)]
      \item \emph{唯一命中性质}：对每个固定的 $s \in A_\kappa$，至多存在一个 $\gamma < \kappa$ 使 $f_{\gamma+1}^\kappa(s) = q_{\gamma+1}(\gamma)$。否则若同时有 $\gamma_1 < \gamma_2 < \kappa$ 满足，用 $f_{\gamma_1+1}^{\gamma_2+1}$ 向下投影并应用投影相容性引理得
        \[
          q_{\gamma_1+1}(\gamma_1) = f_{\gamma_1+1}^\kappa(s) = f_{\gamma_1+1}^{\gamma_2+1}\bigl(q_{\gamma_2+1}(\gamma_2)\bigr) = q_{\gamma_1+1}(\gamma_2).
        \]
        由截断等同判据，$\gamma_1 = \min(\gamma_1+1, \gamma_1) = \min(\gamma_1+1, \gamma_2) = \gamma_1 + 1$（因 $\gamma_2 \ge \gamma_1+1$），矛盾。
      \item 对每个 $s \in S$，记其唯一命中索引为 $\gamma_s < \kappa$（不存在时取 $0$）。因 $|S| < \kappa$ 且 $\kappa$ 是正则基数，上确界 $\alpha := \bigl(\sup_{s \in S} \gamma_s\bigr) + 1 < \kappa$，于是没有任何 $s \in S$ 满足 $f_{\alpha+1}^\kappa(s) = q_{\alpha+1}(\alpha)$。
      \item 然而将 $f_\kappa^{\kappa+1}(\mathsf{L}(S)) = q_\kappa(\kappa)$ 投影到第 $\alpha + 2 < \kappa$ 级：左端是 $f_{\alpha+2}^{\kappa+1}(\mathsf{L}(S)) = \mathsf{L}(f_{\alpha+1}^\kappa[S])$，右端为 $q_{\alpha+2}(\kappa) = \mathsf{L}(\{ q_{\alpha+1}(i) \mid i < \alpha + 2 \})$。故 $f_{\alpha+1}^\kappa[S] = \{ q_{\alpha+1}(i) \mid i < \alpha + 2 \}$。特别地，$q_{\alpha+1}(\alpha) \in f_{\alpha+1}^\kappa[S]$，即存在 $s \in S$ 命中 $\alpha$，矛盾。
    \end{enumerate}
\end{itemize}
因此 $f_\kappa^{\kappa+1}$ 不满射。由稳定的向上传递性，任何使 $G_{\kappa,\theta}$ 稳定的序数 $\delta$ 必满足 $\delta \ge \kappa$。
\end{proof}

\section{第二阶段：下降滤子 $(B_\beta)$、剪枝秩与 $\Stab(G_{\kappa,\theta}) = \kappa + \theta$}

由命题 1.2，$f_\alpha^\beta$ 在全部 $\kappa \le \alpha \le \beta$ 上皆为单射。因此可在 $A_\kappa$ 上定义子对象降链（滤子）：
\[
  B_0 = A_\kappa, \qquad B_\beta := \im\bigl(f_\kappa^{\kappa+\beta} : A_{\kappa+\beta} \to A_\kappa\bigr) \subseteq A_\kappa \qquad (\beta \in \Ord).
\]
对任意子集 $U \subseteq A_\kappa$，记 $H(U) \subseteq G_{\kappa,\theta}(A_\kappa)$ 为受限于 $U$ 的子函子像集：
\begin{itemize}[leftmargin=2em]
  \item 对 $u = \mathsf{L}(S)$，$u \in H(U) \iff S \subseteq U$；
  \item 对 $u = \mathsf{R}(\mu, [a])$，$u \in H(U) \iff \exists b \equiv_{<\mu} a,\; \forall i < \mu,\; b(i) \in U \iff \bigl|\{ i < \mu : a(i) \notin U \}\bigr| < \mu$。
\end{itemize}

\begin{lemma}[$B_\beta$ 的基本性质与递推刻画]
\begin{enumerate}[label=(\arabic*)]
  \item 若 $\beta_1 \le \beta_2$，则 $B_{\beta_2} \subseteq B_{\beta_1}$，且每个 $B_\beta \neq \emptyset$。
  \item \textbf{后继步}：$B_{\beta+1} = f_\kappa^{\kappa+1}\bigl[H(B_\beta)\bigr]$。
  \item \textbf{极限步}：对任何极限序数 $l$，$B_l = \bigcap_{\beta < l} B_\beta$。
  \item \textbf{局部满射等价}：$f_{\kappa+\beta}^{\kappa+\beta+1}$ 为满射当且仅当 $B_{\beta+1} = B_\beta$。
\end{enumerate}
\end{lemma}

\begin{proof}
(1) 由 $f_\kappa^{\kappa+\beta_2} = f_\kappa^{\kappa+\beta_1} \circ f_{\kappa+\beta_1}^{\kappa+\beta_2}$ 与 $A_{\kappa+\beta} \neq \emptyset$ 即得。

(2) 因 $f_{\kappa+1}^{\kappa+\beta+1} = G_{\kappa,\theta}(f_\kappa^{\kappa+\beta})$ 而 $f_\kappa^{\kappa+\beta} : A_{\kappa+\beta} \xrightarrow{\sim} B_\beta$，恰有 $\im(f_{\kappa+1}^{\kappa+\beta+1}) = H(B_\beta)$；再复合 $f_\kappa^{\kappa+1}$ 即得 $B_{\beta+1} = f_\kappa^{\kappa+1}[H(B_\beta)]$。

(3) ($\subseteq$) 由 (1) 显然成立。($\supseteq$) 设 $c \in \bigcap_{\beta < l} B_\beta$，即对每个 $\beta < l$ 存在 $x_\beta \in A_{\kappa+\beta}$ 使 $f_\kappa^{\kappa+\beta}(x_\beta) = c$。对任意 $\beta \le \gamma < l$，
\[
  f_\kappa^{\kappa+\beta}\bigl(f_{\kappa+\beta}^{\kappa+\gamma}(x_\gamma)\bigr) = f_\kappa^{\kappa+\gamma}(x_\gamma) = c = f_\kappa^{\kappa+\beta}(x_\beta),
\]
由尾部单态性（$f_\kappa^{\kappa+\beta}$ 单射）自动推得 $f_{\kappa+\beta}^{\kappa+\gamma}(x_\gamma) = x_\beta$。因此 $(x_\beta)_{\beta < l}$ 满足全部中间层相容性（低层 $\alpha < \kappa$ 由 $x_0 \in A_\kappa$ 往下投影统一补全），从而拼成 $x \in A_{\kappa+l}$ 满足 $f_\kappa^{\kappa+l}(x) = c \in B_l$。

(4) 由 $f_\kappa^{\kappa+\beta}$ 单射，
\[
  f_{\kappa+\beta}^{\kappa+\beta+1} \text{ 满射} \iff \im\bigl(f_{\kappa+\beta}^{\kappa+\beta+1}\bigr) = A_{\kappa+\beta} \iff f_\kappa^{\kappa+\beta}\bigl(\im(f_{\kappa+\beta}^{\kappa+\beta+1})\bigr) = f_\kappa^{\kappa+\beta}(A_{\kappa+\beta}) \iff B_{\beta+1} = B_\beta. \qedhere
\]
\end{proof}

\subsection*{3.1 剪枝秩 $\rho$、延迟算子 $\sigma, \tau_\mu$ 与全部 $< \theta$ 秩的实现}

对 $x \in B_0 = A_\kappa$，若存在序数 $\beta$ 使 $x \in B_\beta \setminus B_{\beta+1}$，记 $\rho(x) = \beta$ 为 $x$ 的\textbf{剪枝秩}。由引理 3.1(4)，存在 $\rho(x) = \beta \iff B_{\beta+1} \subsetneq B_\beta \iff f_{\kappa+\beta}^{\kappa+\beta+1}$ 不满射。

定义两类 $A_\kappa$ 上的延迟算子：
\[
  \sigma(x) := f_\kappa^{\kappa+1}\bigl(\mathsf{L}(\{x\})\bigr), \qquad \tau_\mu(a) := f_\kappa^{\kappa+1}\bigl(\mathsf{R}(\mu, [a])\bigr) \quad (\mu \in \Reg_{<\theta},\; a : \mu \to A_\kappa).
\]

\begin{lemma}[延迟算子的成员判据]
对任意序数 $\beta$：
\begin{enumerate}[label=(\arabic*)]
  \item $\sigma(x) \in B_{\beta+1} \iff x \in B_\beta$；
  \item $\tau_\mu(a) \in B_{\beta+1} \iff \bigl|\{ i < \mu : a(i) \notin B_\beta \}\bigr| < \mu$。
\end{enumerate}
\end{lemma}

\begin{proof}
利用 $B_{\beta+1} = f_\kappa^{\kappa+1}[H(B_\beta)]$ 与 $f_\kappa^{\kappa+1}$ 的单射性，立即等价于 $\mathsf{L}(\{x\}) \in H(B_\beta)$ 与 $\mathsf{R}(\mu, [a]) \in H(B_\beta)$。
\end{proof}

据此立刻得到两个延迟命题与秩实现定理：
\begin{itemize}[leftmargin=2em]
  \item \textbf{后继延迟}：若 $\rho(x) = \alpha$，则 $\rho(\sigma(x)) = \alpha + 1$。
  \item \textbf{极限延迟}：设 $l < \theta$ 为极限序数，记 $\mu := \cof(l)$。因 $l$ 为极限序数且 $l < \theta$，$\mu$ 是无穷正则基数且 $\mu \le |l| < \theta$，即 $\mu \in \Reg_{<\theta}$。取在 $l$ 中严格递增的共尾序列 $(c_j)_{j<\mu}$。若 $\rho(x_{c_j}) = c_j$ 对一切 $j < \mu$ 成立，令 $x_l := \tau_\mu\bigl((x_{c_j})_{j<\mu}\bigr)$，则 $\rho(x_l) = l$：
    \begin{enumerate}[label=(\roman*)]
      \item 对任给 $\beta < l$，由共尾性取 $i_0 < \mu$ 使 $j \ge i_0 \implies c_j \ge \beta \implies x_{c_j} \in B_{c_j} \subseteq B_\beta$；故 $\{j < \mu : x_{c_j} \notin B_\beta\} \subseteq i_0$ 基数 $< \mu$，由成员判据得 $x_l \in B_{\beta+1} \subseteq B_\beta$，从而 $x_l \in \bigcap_{\beta < l} B_\beta = B_l$。
      \item 对一切 $j < \mu$ 均有 $c_j + 1 < l$，故 $B_l \subseteq B_{c_j+1}$，从而 $x_{c_j} \notin B_{c_j+1} \implies x_{c_j} \notin B_l$。于是 $\{j < \mu : x_{c_j} \notin B_l\} = \mu$，基数不小于 $\mu$，故 $x_l \notin B_{l+1}$。
    \end{enumerate}
\end{itemize}

\begin{corollary}[全部 $< \theta$ 秩的实现]
对每一个 $\alpha < \theta$，均存在 $x_\alpha \in B_\alpha \setminus B_{\alpha+1}$。特别地，对任意 $\delta < \kappa + \theta$，$G_{\kappa,\theta}$ 在 $\delta$ 处均不稳定，故 $\Stab(G_{\kappa,\theta}) \ge \kappa + \theta$。
\end{corollary}

\begin{proof}
对 $\alpha < \theta$ 超限递归构造（取 $x_0 = b_\kappa \in B_0 \setminus B_1$，后继步用 $\sigma$，极限步用 $\tau_{\cof(l)}$）即得存在 $x_\alpha$ 使 $\rho(x_\alpha) = \alpha$，从而当 $\kappa \le \delta = \kappa + \alpha < \kappa + \theta$ 时 $f_\delta^{\delta+1}$ 不满射；$\delta < \kappa$ 的情形已由第一阶段排除。
\end{proof}

\subsection*{3.2 $\theta$-正则闭合与精确稳定阶}

\begin{lemma}[小集合下降截断稳定]
设 $S \subseteq B_0 = A_\kappa$ 满足 $|S| < \theta$。则存在 $\beta_* < \theta$ 使得 $S \cap B_{\beta_*} = S \cap B_\theta$。
\end{lemma}

\begin{proof}
对每个 $s \in S \setminus B_\theta$，由 $B_\theta = \bigcap_{\gamma < \theta} B_\gamma$ 选定 $\gamma_s < \theta$ 使 $s \notin B_{\gamma_s}$（对 $s \in S \cap B_\theta$ 置 $\gamma_s = 0$）。由于 $|S| < \theta$ 且 $\theta$ 为正则基数，上确界
\[
  \beta_* := \sup_{s \in S} \gamma_s < \theta.
\]
对任意 $s \in S \cap B_{\beta_*}$，若 $s \notin B_\theta$ 则由 $\gamma_s \le \beta_*$ 得 $B_{\beta_*} \subseteq B_{\gamma_s}$，从而 $s \in B_{\gamma_s}$，矛盾。故 $S \cap B_{\beta_*} = S \cap B_\theta$。
\end{proof}

\begin{theorem}[$\theta$-正则闭合与主定理]
$B_{\theta+1} = B_\theta$，且 $\Stab(G_{\kappa,\theta}) = \kappa + \theta$。
\end{theorem}

\begin{proof}
只需证 $B_\theta \subseteq B_{\theta+1}$。任取 $x \in B_\theta$。由于 $x \in B_\theta \subseteq B_1 = \im(f_\kappa^{\kappa+1})$，存在唯一 $u \in G_{\kappa,\theta}(A_\kappa)$ 使 $x = f_\kappa^{\kappa+1}(u)$；又对一切 $\beta < \theta$ 均有 $x \in B_\theta \subseteq B_{\beta+1} = f_\kappa^{\kappa+1}[H(B_\beta)]$，结合 $f_\kappa^{\kappa+1}$ 的单射性得 $\forall \beta < \theta,\; u \in H(B_\beta)$。分两分支证明 $u \in H(B_\theta)$：
\begin{itemize}[leftmargin=2em]
  \item \textbf{左分支 $u = \mathsf{L}(S)$}：对一切 $\beta < \theta$ 有 $S \subseteq B_\beta$，故 $S \subseteq \bigcap_{\beta < \theta} B_\beta = B_\theta$，从而 $u \in H(B_\theta)$。
  \item \textbf{右分支 $u = \mathsf{R}(\mu, [a])$}（$\mu < \theta$）：令 $S = \im(a) \subseteq A_\kappa$，则 $|S| \le \mu < \theta$。由引理 3.4 取 $\beta_* < \theta$ 使 $S \cap B_{\beta_*} = S \cap B_\theta$。因为 $u \in H(B_{\beta_*})$，我们有 $\bigl|\{i < \mu : a(i) \notin B_{\beta_*}\}\bigr| < \mu$。又对任意 $i < \mu$，$a(i) \in S$，故 $a(i) \in B_\theta \iff a(i) \in B_{\beta_*}$。从而
    \[
      \{i < \mu : a(i) \notin B_\theta\} = \{i < \mu : a(i) \notin B_{\beta_*}\},
    \]
    基数仍严格小于 $\mu$，即 $u \in H(B_\theta)$。
\end{itemize}
由此 $x = f_\kappa^{\kappa+1}(u) \in f_\kappa^{\kappa+1}[H(B_\theta)] = B_{\theta+1}$。

综上 $B_{\theta+1} = B_\theta$，由引理 3.1(4) 与尾部单态性即知 $f_{\kappa+\theta}^{\kappa+\theta+1}$ 为双射；配合对全部 $\delta < \kappa + \theta$ 的否定，精确稳定阶恰好为
\[
  \Stab(G_{\kappa,\theta}) = \kappa + \theta. \qedhere
\]
\end{proof}

\section{特化：$\kappa = \theta = \aleph_1$}

取 $\kappa = \theta = \aleph_1$：此时满足 $\mu < \aleph_1$ 的无穷正则基数仅有 $\mu = \aleph_0$，于是
\[
  G_{\aleph_1, \aleph_1}(X) = \mathcal{P}_{<\omega_1}(X) \;\sqcup\; \bigl(X^\omega / {\equiv_{\mathrm{fin}}}\bigr),
\]
由上述主定理立即给出其精确稳定阶为
\[
  \Stab(G_{\aleph_1, \aleph_1}) = \omega_1 + \omega_1 = \omega_1 \cdot 2.
\]

\end{document}
